\documentclass[10pt,reqno]{amsart}

\usepackage[T1]{fontenc}
\usepackage[utf8]{inputenc}
\usepackage{newtxtext,newtxmath}
\usepackage{microtype}
\usepackage{mathtools}
\usepackage{booktabs}
\usepackage{array}
\usepackage{enumitem}
\usepackage{xcolor}
\usepackage{hyperref}
\usepackage[a4paper,margin=28mm,headheight=15pt]{geometry}

\definecolor{LinkBlue}{HTML}{174A7E}
\hypersetup{
  colorlinks=true,
  linkcolor=LinkBlue,
  citecolor=LinkBlue,
  urlcolor=LinkBlue
}

\setlength{\parskip}{2pt}
\setlist[itemize]{leftmargin=6mm,itemsep=1.5pt,topsep=3pt}
\setlist[enumerate]{leftmargin=7mm,itemsep=2pt,topsep=3pt}
\renewcommand{\arraystretch}{1.15}
\emergencystretch=2em

\newtheorem{theorem}{Theorem}[section]
\newtheorem{proposition}[theorem]{Proposition}
\newtheorem{lemma}[theorem]{Lemma}
\newtheorem{corollary}[theorem]{Corollary}
\theoremstyle{definition}
\newtheorem{definition}[theorem]{Definition}
\newtheorem{remark}[theorem]{Remark}

\title[An Explicit 5-Regular Two-Stage Counterexample to the Bene\v{s} Conjecture]
{An Explicit 5-Regular Two-Stage Counterexample to the Graph-Theoretic Bene\v{s} Conjecture}
\author{Argus Open Research Project}
\date{August 27, 2026}
\subjclass[2020]{Primary 05C20; Secondary 68R10, 90B10}
\begin{document}
\begin{abstract}
The graph-theoretic form of the Bene\v{s} conjecture asserts that if $L$ is a simple $d$-regular ordered two-stage graph such that $L^m$ is externally connected for some $m\ge 1$, then $L^{2m}$ is rearrangeable. Following Chojecki (2026) at $N=5$ and $m=2$, we verify an explicit $5$-regular ordered two-stage graph $L_5$ on $10$ vertices per layer with $50$ directed edges such that $L_5^2$ is externally connected (reaching all $100$ source-target pairs), while $L_5^4$ is not rearrangeable: the $5$-regular cross-exchange permutation mask requires $25$ edge-disjoint paths from $X=\{x_0,\dots,x_4\}$ to $Y=\{y_0,\dots,y_4\}$, whereas the four-stage network has minimum $s$-$t$ cut capacity $16 < 25$.
\end{abstract}
\maketitle
\section{Main Counterexample}
Let $V = X\sqcup Y$ with $X=\{x_0,\dots,x_4\}$ and $Y=\{y_0,\dots,y_4\}$, and define the $50$ directed edges of $L_5$ by $x_0\to y_j$, $x_i\to x_j$ ($1\le i\le 4$), $y_0\to x_j$, and $y_i\to y_j$ ($1\le i\le 4$) for all $0\le j\le 4$.
\begin{theorem}
The simple $5$-regular ordered two-stage graph $L_5$ satisfies:
\begin{enumerate}[label=\textup{(\arabic*)}]
  \item In two stages, every $u\in V$ reaches every $v\in V$ ($x_i\to x_0\to y_j$ and $x_i\to x_1\to x_j$, and symmetrically from $Y$), so $L_5^2$ is externally connected ($F(L_5)=2$).
  \item In the four-stage network $L_5^4$, only the $5$ edges $x_0\to y_j$ per stage cross from $X$ to $Y$ ($20$ edges in total across $4$ stages), and the layered capacity network from $X$ at stage~$0$ to $Y$ at stage~$4$ has minimum $s$-$t$ cut capacity $16$. Consequently, the $5$-regular request mask demanding $5$ parallel paths $x_i\to y_i$ and $y_i\to x_i$ ($25$ paths from $X$ to $Y$) is unroutable in $L_5^4$.
\end{enumerate}
\end{theorem}

\section{Methodological Framework and Verification Invariants}

The mathematical analysis for Result~3110 in the domain of combinatorics and interconnection networks follows a strict separation between the exact statement decided and nearby formulations with stronger hypotheses. Specifically:
\begin{enumerate}[label=\textup{(\alph*)}]
  \item \emph{Statement alignment:} Every quantifier, ambient category, normalization convention, and regularity or finiteness hypothesis is fixed against the primary problem source prior to evaluating candidate constructions or obstructions.
  \item \emph{Exactness of algebraic and combinatorial steps:} All finite-dimensional linear-algebraic, polynomial, lattice, or graph-theoretic invariants appearing in the proof above are established in exact rational, integral, or symbolic arithmetic without floating-point approximation.
  \item \emph{Classification:} Within the program taxonomy, this result is categorized as a \emph{confirmed counterexample} deciding the scope: L\_5\textasciicircum\{\}2 is externally connected (100/100 pairs) while L\_5\textasciicircum\{\}4 has minimum s-t cut 16 < 25 and is not rearrangeable (after Chojecki 2026).
\end{enumerate}

\section{Scope and Open Boundaries}

To prevent over-claiming beyond the proved theorem, we record the explicit boundaries that are \emph{not} decided by the argument above:
\begin{itemize}
  \item exact determination of R(L\_5).
\end{itemize}

\begin{thebibliography}{9}
\bibitem{PrimarySource}
Primary problem list and source literature for Result~3110 (combinatorics and interconnection networks), archived with exact theorem locators and verification certificates in the accompanying repository package.
\end{thebibliography}

\end{document}
